\section*{序論}
10月21日朝までの24時間は、生成AIの応用範囲が音声対話、生命科学、車載体験へ広がる一方、それらを支える推論基盤やエージェント開発環境にも動きがあった。Mistralの音声モデル群、Anthropicのタンパク質設計実験、OpenAIのバイオ分野への投資が目立つ。

インフラ面ではDeepSeekの推論フレームワーク公開、分散推論スタックllm-d、Vercel AI SDKのシェル機能が話題になった。Googleは車載でのパーソナライズ機能を告知し、DeepMindはGemini 4 Argonのデモを披露した。モデルの性能だけでなく、実装・運用までを含めた競争が見える。

本号はGrokの12トピックを収集順に整理する。公式発表と、評価機関・コミュニティによる報告や未検証の主張を区別した。独立した一次資料の再調査は行っていない。
